\section{Related Work}
\label{sec:related-work}

\subsection{Stack Distance and Reuse Distance}

Mattson et al.\ introduced the stack algorithm and the distance between
successive references to a block, and proved that LRU hit behavior at capacity
$C$ is determined by counting distinct blocks referenced since the previous
access \cite{mattson1970}. Hill and Smith developed the miss-ratio curve as a
first-class object, computing all-miss-ratio curves by simulating a set of
capacities and showing how curve shape distinguishes compulsory, capacity and
conflict misses \cite{hill1989}. Both results are exact for the reference
stream they are computed on; neither is a statistical model, and neither
extrapolates to unobserved capacities.

\subsection{Approximate MRC Construction}

Exact stack distance computation costs time linear in the trace length times
the number of distinct blocks. Counter stacks replace per-reference search with
a compact data structure that yields stack distance histograms for many
capacities in one pass \cite{wires2014}. SHARDS obtains MRCs at scale by
sampling references into miniature simulations, correcting the resulting
downsampling bias analytically, and aggregating per-simulation histograms
\cite{waldspurger2015}. These methods are the ground truth against which a
model is judged here: simulation is accurate up to sampling error, and its cost
is what a model is meant to avoid on every capacity query.

\subsection{Reuse-Distance Analysis}

Reuse distance analysis extends the single-parameter stack model to arbitrary
replacement policies by tracking the reuse distance of each request and
replaying it against a policy-specific recency structure
\cite{beckmann2016,beckmann2017}. This line of work establishes that reuse
distance is a sufficient statistic for a broad class of policies, which is the
premise the present model relies on for LRU and only for LRU. Production trace
studies supply the workload evidence: the Twitter cache analysis reports
reuse-distance and miss-ratio structure across hundreds of in-memory clusters
and motivates cluster-level heterogeneity as a modelling concern
\cite{yang2020}.

\subsection{Survival Analysis in Systems Contexts}

Reuse distance is a positive, integer-valued waiting time, observed under a
truncation that depends on the observation window. That is a censored
time-to-event problem, and the discrete hazard formulation used here follows
the Kaplan-Meier product-limit estimator for right-censored data
\cite{kaplan1958} and the proportional hazards regression of Cox
\cite{cox1972}. Grouped and interval-censored observations, which arise when
reuse distance is binned, are treated by the Turnbull estimator
\cite{turnbull1976}. Gradient boosting supplies the flexible hazard component:
gradient-boosted decision trees \cite{friedman2001} as implemented in LightGBM
\cite{ke2017}, with hyperparameters selected by Optuna \cite{akiba2019}. The
combination of a nonparametric hazard with a tree ensemble is standard in
survival regression; the contribution here is the mapping of that hazard onto
cache capacity rather than onto a clinical or reliability endpoint.
